\documentclass[11pt]{article}
\usepackage[margin=2.6cm]{geometry}
\usepackage{amsmath,amssymb,amsthm}
\usepackage{xcolor}
\newcommand{\seg}[1]{\textcolor{gray}{\texttt{\small #1}}}
\newcommand{\Hom}{\mathrm{Hom}}
\newcommand{\I}{\mathbb{I}}
\newcommand{\Ext}{\mathrm{Ext}}
\newcommand{\ext}{\mathrm{ext}}
\newcommand{\glue}{\mathrm{glue}}
\newcommand{\Fib}{\mathrm{Fib}}
\newcommand{\full}{\mathrm{full}}
\theoremstyle{definition}
\newtheorem{definition}{Definition}
\newtheorem{lemma}[definition]{Lemma}
\newtheorem{construction}[definition]{Construction}
\theoremstyle{remark}
\newtheorem*{remark}{Remark}

\title{The filling step of Theorem 7:\\ $\ext(\alpha)$ is an equivalence\\[2pt]
\large the mathematics of \texttt{seg/model/GExtEquiv2.seg}}
\date{15 October 2026}
\author{}

\begin{document}
\maketitle

\noindent
This note states, as a mathematical argument, what the file
\texttt{GExtEquiv2.seg} proves. It follows the cleaned definition, in which
a cofibration is a single decidable sieve rather than a list of generators;
the file is being brought in line with it. Each paragraph corresponds to a group of
declarations of the file; their names are given in grey. A declaration
marked \texttt{opaque} in the file is a lemma here: after it is stated,
only its statement is used.

\section*{Setting}

The base is a small category $\mathcal C$ with an object $[1]$ and, for
every object $X$, two maps $e_0,e_1:X\to[1]$, natural in $X$ (the constant
maps at the end points). We write $\I=y[1]$ for the presheaf represented by
$[1]$, so that $\I(W)=\Hom(W,[1])$. We do not assume that $\mathcal C$ has
products: posets
and $\Delta$ are both instances. For each object $X$ there is a set
$\Phi(X)$ of decidable sieves on $X$, stable under precomposition. A map
$f:V\to X$ \emph{lands in} $Q\in\Phi(X)$ if $f\in Q$.

\begin{itemize}
\item $G$ is a presheaf (the context) and $A$ is a type over $G$. We write
  $A(X,\rho)$ for its fibre at $\rho\in G(X)$, and $E_A$ for the total
  presheaf of the pairs $(\rho,a)$.
\item A \emph{cofibration} $\psi$ over $G$ is a decidable sieve
  $\psi_\rho\in\Phi(X)$ for each $\rho\in G(X)$, natural in $\rho$:
  $\psi_{\rho f}=f^*\psi_\rho$. The sets $\Phi(X)$ are closed under finite
  unions. The point $\rho$ is \emph{full}, written $\full(X,\rho)$, if
  $\mathrm{id}_X\in\psi_\rho$. By naturality, a map $f:W\to X$ is full
  (meaning $\rho f$ is full) exactly when $f\in\psi_\rho$. Being full is a
  proposition.
\item $G|\psi$ is the context of the full points. $B$ is a type over
  $G|\psi$, and $\alpha:B\to A|\psi$ is a map over $G|\psi$.
\end{itemize}

\paragraph{Lines.}
A \emph{line in $A$ over $\rho\in G(X)$} is a family
$\ell(W,f,j)\in E_A(W)$, for $f:W\to X$ and $j\in\I(W)$, which lies over
$\rho f$ and is natural:
\[
\ell(W,f,j)\cdot g=\ell(V,fg,jg)\qquad(g:V\to W).
\]
It is thus a map $yX\times\I\to E_A$ over $\rho$, which is how the
family path object is defined without products in $\mathcal C$. Its ends are
$\ell(e_0)=\ell(X,\mathrm{id},e_0)$ and
$\ell(e_1)=\ell(X,\mathrm{id},e_1)$. This is the family path object
(\seg{GFam.famPath}). The \emph{constant line} at $\rho$ is the line of the
base $G$ given by $(W,f,j)\mapsto\rho f$ (\seg{constF}).

\paragraph{Fibres and equivalences.}
Let $m:B\to A$ be a map over a context $D$. For $w\in E_A(X)$, an element of
the \emph{fibre} $\Fib_m(X,w)$ is a pair $(b,\ell)$. Here $b\in B$ lies over
the point of $w$, and $\ell$ is a line in $A$ over that point with
$\ell(e_0)=m(b)$ and $\ell(e_1)=w$ (\seg{GEquiv.HFib}: \seg{fibB},
\seg{fibL}, \seg{fibL-0}, \seg{fibL-1}).

An \emph{equivalence structure} on $m$ (\seg{EquivOp}) extends every partial
element of a fibre to a total one. The extension agrees with the sides
(\seg{kside}), is uniform (\seg{kunif}), and is the value of any side that
contains the identity (\seg{kfull}).

\paragraph{The data.}
The file works with the following data.
\begin{itemize}
\item An equivalence structure $e_\alpha$ on $\alpha$ (\seg{eal}).
\item A filling structure $k_A$ for $A$ at $e_1$ (\seg{kA : FibAt e1}). It
  extends a partial line in $A$, over a line of $G$, whose end at $e_1$ is
  given totally, to a total line with the same end. Uniformity and
  agreement with the sides are as above.
\end{itemize}

\paragraph{$\Ext(\alpha)$.}
Let $\rho\in G(X)$. If $\rho$ is full, $\Ext(\alpha)(X,\rho)=B(X,\rho)$.
Otherwise it is the set of pairs $(a,b)$ (\seg{GExt}, \seg{Pairs}) where:
\begin{itemize}
\item $a\in A(X,\rho)$;
\item $b$ assigns to each full $f:W\to X$ an element $b_f\in B(W,\rho f)$;
\item $b$ is natural: $b_f\cdot\chi=b_{f\chi}$;
\item $\alpha(b_f)=a\cdot f$.
\end{itemize}
The map $\glue$ sends such a pair to $\Ext(\alpha)(X,\rho)$ in both cases.
If $\rho$ is full, it takes $b_{\mathrm{id}}$. The map
$\ext:\Ext(\alpha)\to A$ satisfies $\ext(\glue(a,b))=a$
(\seg{glue-ext}). Because $\glue$ handles both cases, the argument below
never needs to ask whether $\rho$ is full.

\section*{The problem}

Fix $X$ and $w=(\rho,a)\in E_A(X)$. Let $p$ be a partial element of
$\Fib_{\ext}(X,w)$. It has sides $Q'_0,\dots,Q'_{n-1}\in\Phi(X)$ and
values $p_l(V,f)\in\Fib_{\ext}(V,wf)$ for $f$ landing in $Q'_l$. These
values are natural in $f$ and agree on overlaps.

We construct an element $v_0\in\Fib_{\ext}(X,w)$ (\seg{v0}). In
\texttt{GExtEquiv3} it is shown that $v_0$ extends $p$ and is uniform; this
proves that $\ext(\alpha)$ is an equivalence (half of Theorem 7). The idea
is a single filling in $A$, over the constant line at $\rho$. One family of
sides comes from $p$. The other comes from extending, with $e_\alpha$, what
$p$ says at the full maps.

\section*{The argument}

\begin{construction}[$\theta$, \seg{th}]
Let $f:W\to X$ be full. Restrict $p$ along $f$. At a full point, a fibre
element of $\ext$ is a fibre element of $\alpha$: on the $B$ side,
$\Ext(\alpha)$ is $B$ there (\seg{toFA}, \seg{Tr}). This gives a partial
element of $\Fib_\alpha(W,wf)$. Let
\[
\theta_f\in\Fib_\alpha(W,wf)
\]
be its extension by $e_\alpha$ (\seg{kE}).
\end{construction}

\begin{lemma}[\seg{th-nat}, \seg{th-full}, \seg{th-cong}]\label{lem:theta}
\leavevmode
\begin{enumerate}
\item $\theta$ is natural: $\theta_f\cdot\chi=\theta_{f\chi}$ for
  $\chi:W'\to W$.
\item If $f$ lands in $Q'_l$, then $\theta_f$ is $p_l(W,f)$ read as a fibre
  element of $\alpha$.
\item $\theta_f$ does not depend on the proof that $f$ is full.
\end{enumerate}
\end{lemma}

\begin{proof}
\begin{enumerate}
\item Restricting the translated partial element along $\chi$ gives the
  translation of the restricted one. The claim then follows from the
  uniformity of $e_\alpha$ (\seg{kE-act}).
\item Along $f$, the side $Q'_l$ contains the identity of $W$
  (\seg{lands-idr}). An extension is the value of a side that contains the
  identity (\seg{kE-full}). Restricting along the identity changes nothing
  (\seg{pv-hom}).
\item Being full is a proposition.
\end{enumerate}
\end{proof}

\begin{construction}[the base point, \seg{x0}]
The line of $G$ is the constant line at $\rho$. Its prescribed total end at
$e_1$ is $w$ (\seg{x0-E}). Its restriction along $f$ ends at $wf$
(\seg{x0-act}).
\end{construction}

\begin{construction}[the sides]
\leavevmode
\begin{enumerate}
\item[(S1)] For each side $Q'_l$ of $p$, take the line of $p_l(V,f)$
  (\seg{S1}). It lies over $\rho f$, ends at $wf$ at $e_1$ by
  \seg{fibL-1}, and is natural in $f$.
\item[(S2)] On the sieve $\psi_\rho$, take the image in $A$ of the line
  of $\theta_f$ (\seg{S2}). It is defined because every $f\in\psi_\rho$ is
  full. The line lies over $\rho f$ and ends at $wf$. It is natural in $f$
  by Lemma~\ref{lem:theta}(1).
\end{enumerate}
\end{construction}

\begin{lemma}[compatibility: \seg{c11}, \seg{c22}, \seg{c12}]
The sides (S1) and (S2) agree wherever two of them are both defined.
\end{lemma}

\begin{proof}
Two elements over the same point, whose lines take the same values, are
equal (\seg{tsEq}). So it is enough to compare lines.
\begin{itemize}
\item Two sides (S1): this is the compatibility of $p$.
\item The side (S2) with itself: this is Lemma~\ref{lem:theta}(3).
\item One side of each kind: $f$ lands in $Q'_l$ and in $\psi_\rho$. By
  Lemma~\ref{lem:theta}(2), $\theta_f$ is $p_l(V,f)$, and the line of a
  translated fibre element is the line it came from (\seg{toFA-L}).
\end{itemize}
\end{proof}

The sides are joined into one partial line $J$ with $n+1$ sides
(\seg{GCofib.Join}).

\begin{construction}[the filler, \seg{fill}]
Let $F=k_A(x_0,J)$. This is a line in $A$ over the constant line at $\rho$.
Its values are $F(W,f,j)\in A(W,\rho f)$ (\seg{Fv}, \seg{Fn}, \seg{Fpt}).
\end{construction}

\begin{lemma}[\seg{F1}, \seg{F-S1}, \seg{F-S2}, \seg{F0-act}]\label{lem:F}
\leavevmode
\begin{enumerate}
\item $F(e_1)=w$.
\item On $Q'_l$, $F$ is the line of $p_l$.
\item On $\psi_\rho$, $F$ is the line of $\theta$.
\item $F(e_0)\cdot f=F(W,f,e_0)$.
\end{enumerate}
\end{lemma}

\begin{proof}
Items 1 to 3 hold because $k_A$ keeps the given end and agrees with the
sides (\seg{Fside}), and because the joined partial line restricts to the
sides of $p$ and to the side $\psi_\rho$ (\seg{J.atL}, \seg{J.atR}). Item 4 is the naturality of
$F$.
\end{proof}

\begin{construction}[the pair, \seg{u}]
Let $a_0=F(e_0)\in A(X,\rho)$ (\seg{a0}). For full $f$, let $b_f$ be the
$B$ component of $\theta_f$ (\seg{bE}). Then $u=(a_0,b)$ is a pair
(\seg{mkPairs}):
\begin{itemize}
\item $b$ is natural by Lemma~\ref{lem:theta}(1) (\seg{pnat0});
\item $\alpha(b_f)=a_0\cdot f$ (\seg{pal0}) by the chain
\[
\alpha(b_f)\overset{\text{fibre}}{=}\text{line of }\theta_f\text{ at }e_0
\overset{\text{L.\ref{lem:F}(3)}}{=}F(W,f,e_0)
\overset{\text{L.\ref{lem:F}(4)}}{=}a_0\cdot f,
\]
where Lemma~\ref{lem:F}(3) applies because the full map $f$ lies in
$\psi_\rho$.
\end{itemize}
\end{construction}

\begin{construction}[the fibre element, \seg{v0}]
Let $v_0=(\glue(u),F)$. It lies in $\Fib_{\ext}(X,w)$. Indeed $F$ is a line
over $\rho$, with
\[
F(e_0)=a_0=\ext(\glue(u))\quad(\seg{u-ext},\ \seg{F0})
\qquad\text{and}\qquad
F(e_1)=w\quad(\text{Lemma~\ref{lem:F}(1)}).
\]
\end{construction}

\begin{remark}
The box proof (\texttt{ExtEquiv2}, \texttt{ExtKit}) filled a box over the
cylinder $X\times[1]$, an object of $\mathcal C$ for posets. Here a single partial \emph{line} over the constant
line is filled. So the argument uses only the presheaf $yX\times\I$, never a product in
$\mathcal C$, and it applies
both to posets and to $\Delta$. No $\forall$ along the interval is needed,
because the line of $G$ is constant: the side (S2) is $\psi_\rho$
itself. That quantifier enters only in the
composition part of Theorem 7 (step~5).
\end{remark}

\section*{Correspondence}

\begin{center}
\begin{tabular}{ll}
mathematics & \texttt{GExtEquiv2.seg}\\\hline
$\theta_f$; Lemma~\ref{lem:theta} & \texttt{th}; \texttt{th-nat}, \texttt{th-full}, \texttt{th-cong}\\
base point (constant line, end $w$) & \texttt{cL}, \texttt{x0}, \texttt{x0-E}, \texttt{x0-act}\\
sides (S1), (S2) & \texttt{S1[l].SO}, \texttt{S2.SO} (\texttt{GMapLine.LineSide})\\
compatibility & \texttt{tsEq}, \texttt{c11}, \texttt{c22}, \texttt{c12}; \texttt{J}\\
filler $F$; Lemma~\ref{lem:F} & \texttt{fill}, \texttt{Fv}; \texttt{F1}, \texttt{Fside}, \texttt{F-S1}, \texttt{F-S2}, \texttt{F0-act}\\
the pair $u$ & \texttt{a0}, \texttt{bE}, \texttt{pal0}, \texttt{pnat0}, \texttt{u}\\
$v_0$ & \texttt{u-ext}, \texttt{gb}, \texttt{F0}, \texttt{v0}
\end{tabular}
\end{center}

\end{document}
